\documentclass[11pt,a4paper]{article}

\usepackage[utf8]{inputenc}
\usepackage[T1]{fontenc}
\usepackage[spanish,es-noshorthands]{babel}
\usepackage[margin=2.8cm]{geometry}
\usepackage{titlesec}
\usepackage{enumitem}
\usepackage{tabularx}
\usepackage{array}
\usepackage{booktabs}
\usepackage{xcolor}
\usepackage[colorlinks=true,linkcolor=black,urlcolor=blue,citecolor=black]{hyperref}
\usepackage{parskip}
\usepackage{microtype}

\newcommand{\code}[1]{\texttt{\small #1}}
\emergencystretch=3em
\sloppy

\titleformat{\section}{\Large\bfseries}{\thesection}{0.8em}{}
\titleformat{\subsection}{\large\bfseries}{\thesubsection}{0.8em}{}
\titlespacing*{\section}{0pt}{1.6em}{0.8em}
\setlist[itemize]{leftmargin=1.6em}

\begin{document}

% ---------- Portada ----------
\begin{titlepage}
    \centering
    \vspace*{5cm}
    {\Huge\bfseries DevLab Interface Library\par}
    \vspace{0.5em}
    {\Large\bfseries Documentación Técnica\par}
    \vspace{2em}
    {\large 23 de septiembre de 2026 \,\textperiodcentered\, Desarrollador: Jonathan Mejorado\par}
    \vfill
    {\normalsize UNIT Electronics --- DevLab Ecosystem\par}
    \vspace{0.6em}
    {\normalsize \url{https://github.com/UNIT-Electronics-MX/unit_devlab_interface_library}\par}
    \vspace{2cm}
\end{titlepage}

% ---------- Cuerpo ----------

\section*{Resumen}

\textbf{DevLab Interface} es la capa de acceso a buses del ecosistema DevLab de UNIT Electronics: una biblioteca de Arduino/C++ que define un único contrato de acceso a registros (lectura/escritura de 8, 16 y 32 bits, más helpers de bit-field) y permite que cualquier transporte físico --- I2C hoy, SPI en desarrollo --- lo implemente con solo dos primitivas crudas (\code{writeBytes}, \code{readBytes}).

La idea central: \textbf{un driver de sensor se escribe una sola vez contra el contrato, y corre sin cambios sobre cualquier bus que lo implemente.} Hoy eso cubre I2C de 7 y 10 bits, con recuperación automática de bus ante fallos, y soporte probado en ESP32, ESP32-C6/C5/H2, RP2040 y STM32. SPI está planeado como el siguiente transporte concreto.

En versión \textbf{1.1.2}, la biblioteca ya es la base de referencia para integrar sensores externos y dispositivos propios de DevLab (como el protocolo DDP de direccionamiento dinámico) sin acoplar cada driver a un bus o a una plataforma específica. Este documento explica la arquitectura, el estado actual y hacia dónde va --- no es un manual función por función; esa referencia vive en el código y en los ejemplos del repositorio.

\section*{El problema que resuelve}

El ecosistema DevLab integra sensores propios y de terceros (Bosch BMI323, MPU6050, etc.) sobre múltiples microcontroladores (ESP32, RP2040, STM32). Sin una capa común, cada combinación sensor \(\times\) bus \(\times\) plataforma tiende a producir su propio código de acceso a registros: distintas formas de escribir un puntero de registro, hacer un repeated-start, descartar bytes de relleno que algunos chips anteponen a sus lecturas (el BMI323, por ejemplo, antepone bytes muertos antes del dato real), o inicializar los pines SDA/SCL, que cada core de Arduino resuelve de forma distinta (ESP32, RP2040 y STM32duino tienen firmas de \code{begin()} diferentes).

Esa fragmentación tiene un costo directo para el equipo: un driver de sensor escrito para I2C no se reutiliza en SPI sin reescribir su lógica de acceso, un bug de inicialización de bus se corrige en una plataforma y reaparece en otra, y cada integración nueva reinventa el manejo de casos borde (direcciones de 10 bits, buses que quedan bloqueados tras un reset, dispositivos que necesitan re-direccionarse en tiempo de ejecución).

DevLab Interface ataca esto separando dos cosas que solían mezclarse: \textbf{el contrato de acceso a registros} (qué significa leer/escribir un registro de N bits) y \textbf{el transporte físico} que lo implementa. El driver de un sensor se programa una sola vez contra el contrato; el transporte concreto se elige por bus, no por sensor.

\section*{Arquitectura del contrato: \code{DevLab\_BusIO}}

El corazón de la biblioteca es una sola clase plantilla, \code{DevLab\_BusIO<Derived>}, resuelta mediante \textbf{CRTP} (Curiously Recurring Template Pattern) en tiempo de compilación. La decisión de diseño es deliberada: en lugar de una interfaz virtual (con vtable y despacho en tiempo de ejecución), cada bus concreto se declara como \code{class DevLab\_I2C\_Interface : public DevLab\_BusIO<DevLab\_I2C\_Interface>}, y el compilador resuelve las llamadas de forma directa. En microcontroladores con RAM y ciclos limitados, esto elimina el costo de indirección sin renunciar a la abstracción.

Un bus concreto solo tiene que implementar dos primitivas:

\begin{itemize}
    \item \code{writeBytes(reg, data, len)}
    \item \code{readBytes(reg, data, len)}
\end{itemize}

Sobre esas dos, \code{DevLab\_BusIO} construye --- una sola vez, compartido por todos los buses --- el resto del contrato: \code{readRegister8/16/32}, \code{writeRegister8/16/32} (little-endian, como la mayoría de los layouts de registros de sensores) y los helpers de bit-field \code{readBits}, \code{readBits16}, \code{readBits32}. Un driver de sensor escrito contra este contrato funciona sin cambios sobre I2C hoy y sobre SPI el día que ese transporte esté disponible.

Otra decisión relevante: \textbf{el contrato es header-only, pero los buses concretos no lo son.} \code{DevLab\_I2C\_Interface}, \code{DevLab\_I2C10\_Interface} y el orquestador son pares \code{.h}/\code{.cpp} normales, no plantillas. Esto mantiene los binarios de firmware pequeños y los tiempos de compilación razonables, algo que importa en proyectos con múltiples sensores y múltiples archivos fuente.

\section*{La capa I2C actual}

I2C es hoy el único transporte implementado, pero está resuelto en varias capas para no repetir código entre variantes:

\textbf{\code{DevLab\_I2C\_Common}} concentra la plomería compartida por todas las variantes de I2C: el puntero a \code{TwoWire}, el reloj, el conteo de bytes de relleno (``dummy bytes'') a descartar en cada lectura, y la lógica de \code{begin()}. Es herencia simple, sin virtuales --- mismo criterio de costo cero que el contrato ---, y evita el copy-paste entre las variantes de 7 y 10 bits.

\textbf{\code{DevLab\_I2C\_Interface}} implementa el contrato para direccionamiento I2C estándar de 7 bits: escribe el puntero de registro, hace repeated-start y lee N bytes. Los bytes de relleno que algunos chips anteponen a sus lecturas (el caso de referencia es el Bosch BMI323) se declaran una vez en el constructor y cada helper los descarta automáticamente --- el driver del sensor nunca calcula offsets de bytes a mano.

\textbf{\code{DevLab\_I2C10\_Interface}} implementa el mismo contrato para direccionamiento de 10 bits (el esquema reservado \code{0b11110XX} de la especificación I2C, con las dos primeras líneas de dirección en la cabecera y los 8 bits restantes en un segundo byte). Se mantiene como clase separada --- en vez de un flag en tiempo de ejecución dentro de \code{DevLab\_I2C\_Interface} --- para que las transacciones de 7 bits, que son la inmensa mayoría, nunca paguen el costo de una verificación que no necesitan.

\textbf{Recuperación de bus}, agregada recientemente: si un dispositivo queda a mitad de una transacción (por ejemplo, tras un reset), SDA puede quedar retenido en bajo y bloquear el bus para todos los demás dispositivos. \code{DevLab\_I2C\_Common} implementa la secuencia de recuperación estándar (hasta 18 pulsos de reloj manuales en SCL para liberar al esclavo, seguidos de una condición STOP), expuesta como \code{devlabBeginI2CBusRecovered()}. El orquestador ya la usa a través de \code{beginRecovered()}; es candidata natural para exponerse también en \code{DevLab\_I2C\_Interface}/\code{DevLab\_I2C10\_Interface} si el equipo lo necesita en dispositivos de dirección fija.

\section*{El orquestador I2C y el protocolo DevLab DDP}

\code{DevLab\_I2C\_Interface} y \code{DevLab\_I2C10\_Interface} asumen un dispositivo de dirección fija --- el modelo típico de un sensor. Pero hay una clase de dispositivos DevLab que necesita otra cosa: escanear el bus, identificar qué hay conectado, y re-direccionar dispositivos en tiempo de ejecución. Para eso existe \textbf{\code{DevLab\_I2C\_Orchestrator}}, pensado explícitamente para protocolos a nivel de bus como \textbf{DevLab DDP.}

Su modelo de transacción es distinto al de una interfaz de registros: en vez de ``escribir registro + repeated-start + leer'', sigue un ciclo de \textbf{comando/respuesta} --- escribe un byte de comando (con STOP), deja que el dispositivo tome su tiempo de procesamiento, y luego emite una lectura nueva. El orquestador no asume una respuesta inmediata; el dispositivo decide cuándo está listo. Expone tres operaciones: \code{ping()} (verificar si una dirección responde), \code{writeCommand()} (comando sin retorno) y \code{transact()} (comando con retorno, con retardo configurable).

A diferencia de las interfaces de registro, una sola instancia del orquestador habla con \textbf{cualquier dirección del bus} --- no guarda una dirección propia ---, lo que lo hace la pieza natural para escanear el bus buscando dispositivos DevLab, identificarlos y asignarles dirección en frío. Esto es lo que habilita, a nivel de infraestructura, un modelo de \textbf{dispositivos plug-and-play}: un módulo DevLab puede conectarse al bus sin dirección fija preasignada de fábrica y ser descubierto por el host en tiempo de ejecución.

El orquestador es también el primer consumidor de la recuperación de bus (\code{beginRecovered()}), coherente con su rol: si va a escanear un bus compartido por varios dispositivos, es el punto donde un esclavo atascado tiene más probabilidad de aparecer.

\section*{Compatibilidad de plataformas}

La inicialización de pines SDA/SCL es uno de los puntos donde más diverge cada core de Arduino, y es exactamente el tipo de detalle que la biblioteca absorbe en \code{DevLab\_I2C\_Common::beginCommon()} para que el resto del código no tenga que pensarlo:

\begin{center}
\begin{tabularx}{\textwidth}{@{} >{\raggedright\arraybackslash}p{4.1cm} >{\raggedright\arraybackslash}p{1.9cm} X @{}}
\toprule
\textbf{Plataforma} & \textbf{Estado} & \textbf{Notas} \\
\midrule
ESP32 (incl. C6 / C5 / H2) & Probado & \code{Wire.begin(sda, scl, clock)} en una sola llamada \\[0.4em]
RP2040 (arduino-pico) & Probado & Pines se fijan antes de \code{begin()}, sin argumentos \\[0.4em]
STM32 (STM32duino core) & Probado & \code{setSDA()}/\code{setSCL()} antes de \code{begin()}; incorporado en la versión 1.1.2 \\[0.4em]
Otras placas Arduino-compatibles con AVR & Probado & Usa \code{Wire.begin()} estándar sin pines explícitos \\
\bottomrule
\end{tabularx}
\end{center}

En la práctica, el trabajo de mantenimiento recurrente de esta capa: cada nuevo core de Arduino que el equipo quiera soportar entra aquí, una sola vez, y queda disponible para todos los buses y todos los drivers que se construyan sobre el contrato.

\section*{En desarrollo y hoja de ruta}

\textbf{SPI concreto (\code{DevLab\_SPI\_Interface}) --- planeado, es el siguiente hito.} El contrato ya está diseñado para ser agnóstico de transporte, falta la clase concreta que lo formalice dentro de este repositorio, en paralelo a como hoy existen \code{DevLab\_I2C\_Interface} y \code{DevLab\_I2C10\_Interface}.

\textbf{Recuperación de bus fuera del orquestador.} \code{devlabBeginI2CBusRecovered()} ya vive en \code{DevLab\_I2C\_Common} y hoy solo se expone públicamente a través de \code{DevLab\_I2C\_Orchestrator::beginRecovered()}. Extender ese mismo \code{beginRecovered()} a \code{DevLab\_I2C\_Interface} y \code{DevLab\_I2C10\_Interface} es un cambio pequeño y ya soportado por la base común --- relevante para cualquier sensor de dirección fija que pueda quedar en un bus compartido con dispositivos DDP.

\section*{Para qué servirá: casos de uso concretos}

\textbf{Portar un driver de sensor entre buses sin reescribirlo.} Un driver contra \code{DevLab\_BusIO} (como \code{DevLab\_BMI323}, ya usado en los ejemplos de referencia) se valida hoy sobre I2C y se reutiliza sobre SPI en cuanto exista \code{DevLab\_SPI\_Interface}, sin tocar la lógica del sensor. Para el equipo, esto significa que la elección de bus en una placa nueva deja de ser una decisión que obliga a reescribir drivers.

\textbf{Dispositivos DevLab plug-and-play (DDP).} El orquestador es la base para que módulos DevLab se descubran, identifiquen y direccionen en el bus sin configuración manual de direcciones --- el tipo de infraestructura que necesita cualquier producto que quiera soportar ``conecta y listo'' entre múltiples módulos en el mismo bus I2C.

\textbf{Integración rápida de sensores de terceros.} El manejo de bytes de relleno por dispositivo (declarado una vez en el constructor) y el \code{getWhoAmI()} de validación genérica bajan el costo de traer un sensor externo nuevo al ecosistema: no hay que escribir código de bus desde cero, solo mapear sus registros sobre el contrato existente.

\textbf{Firmware resiliente en campo.} La recuperación de bus I2C importa en cualquier producto donde el dispositivo pueda reiniciarse o desconectarse sin que el host lo controle (un módulo que se hot-plug, un sensor que se resetea por ruido eléctrico): sin ella, un bus atascado requiere ciclar la alimentación completa del sistema para recuperarse.

\textbf{Portabilidad entre líneas de producto con distintos MCU.} Con ESP32, RP2040, STM32 y AVR ya cubiertos, un mismo driver de sensor corre sin cambios en productos que usan distintas familias de microcontrolador.

\end{document}
